\section{Operating System} \label{app:os}
\subsection{Dataset details}
\vpara{Construction Details.}
Each evaluation sample in OS dataset encompasses following contents:

\begin{itemize}[leftmargin=1.5em,itemsep=0pt,parsep=0.2em,topsep=0.0em,partopsep=0.0em]
    \item \textbf{Instruction.} The description of the problem in natural language that needs LLMs to solve.
    \item \textbf{Docker Environment.} The starting up docker image (e.g., preset default \texttt{local-os/default}).
    \item \textbf{Initialization Script (Optional).} The bash scripts that need to be executed independently (\texttt{docker exec}) before the interaction starts (e.g., user configurations, files, system statuses).
    \item \textbf{Start Script (Optional).} The bash scripts executed after shell is created and before interaction.
    \item \textbf{Checking Pipeline.} The checking method to judge the correctness of LLMs answer or operation.
    \item \textbf{Example Script (Optional).} The bash scripts that serve as reference solutions. In other words, if executing them in the interaction, results are correct. Only for unit tests that introduced below.
\end{itemize}



We design two types of tasks in the OS evaluation beyond conventional QA-only evaluation.
\begin{itemize}[leftmargin=1.5em,itemsep=0pt,parsep=0.2em,topsep=0.0em,partopsep=0.0em]
    \item \textbf{Question Answering (QA)}: LLMs need to output commands to solve specific questions in OS (e.g., aggregate numbers, view file contents). In this case, they must commit answers finally.
    \item \textbf{Operation}: LLMs need to output commands to do some verifiable operations on the operating system (e.g., change file/user states). In this case, they do not need to commit final answers. 
\end{itemize}
Thanks to the checking pipeline, two types of tasks can be evaluated in a unified solution.

Collecting challenging queries regarding OS could be difficult.
In practice, about half of our instructions are created or collected from humans, while the other half are mostly QA problems generated by \texttt{gpt-4} and strictly filtered by passing the unit tests (i.e., yield correct answers/states).

For human instructions, we first gather 6000 real problems and solutions with \texttt{bash} or \texttt{shell} tag from Stack Overflow\footnote{https://stackoverflow.com/}. Then we sort them by the score (count of likes). We invite 8 annotators majored in programming to select challenging ones. 
For each selected problem, they create one or more task instructions and write a detailed problem description, the initialization script, the starting script, and the checking pipeline. 
Finally, we conduct a cross verification for each evaluation sample to make sure it's correct. For each problem, it takes about 2 hours to do the annotation.


For generated problems, our unit test contains the following parts. 
1) Initialization Script Correction:
we execute the initialization script and remove samples with wrong initialization whose exit code does not equal to $0$.
2) Example Code Correction: we execute the example code and the checking pipeline to judge the correctness of the answer. We remove samples with wrong answers.

In the end, we curate 144 high-quality diverse OS evaluation samples accompanied with testing interactive environments and corresponding checking pipelines (i.e., scripts).
Agents are prompted with 1-shot CoT to better format their responses (Cf. Appendix~\ref{app:os}).

\vpara{Evaluation Setup.}
For each problem (i.e., instruction), the execution can be divided into 3 parts. 

\begin{itemize}[leftmargin=1.5em,itemsep=0pt,parsep=0.2em,topsep=0.0em,partopsep=0.0em]
    \item \textbf{Initialization.} We create a docker container with a specific image, and we run an initialization bash script to set up environments specified by the instruction.
    \item \textbf{Interaction.} We start a new shell in this docker, and run the starting bash script specified by the instruction. 
    Then the LLM to test is fed with a piece of instruction and the problem description. 
    It starts interaction with the shell. 
    In each round, two actions are provides. 
    One is to run bash script, which allows the model to generate and run a series of commands in the shell. 
    The other is to commit answer, which allows the model to terminate the interaction process. 
    It's notable that the model will be judged that it fail to solve the problem if exceeding round limit (8 by default).
    \item \textbf{Checking.} For each problem, there is a checking pipeline containing a list of scripts $f_1, f_2, \cdots, f_n$, where $f_k$ denotes the $k$-th script piece in the pipeline. 
    For $f_k$, the answer of the model, $o_0$, and the output of $f_t(t< k)$, $o_t$, will be fed as input arguments into $f_k$, i.e., $o_k=f_k(o_0, o_1, \cdots, o_{k-1})$. 
    The result is correct if and only if all the scripts exit with code $0$.
\end{itemize}

\vpara{Metrics.}
We measure the \textbf{Success Rate} for LLMs to solve problems in the execution. There are only two final status for each item of the problems, wrong or correct.


\subsection{Actions}
In OS evaluation, we design two major types of actions: \textit{bash} and \textit{commit}.

\begin{itemize}[leftmargin=1.5em,itemsep=0pt,parsep=0.2em,topsep=0.0em,partopsep=0.0em]
    \item \textbf{Bash}: which launches a bash command (using textual inputs in \textit{content} field)
    \item \textbf{Commit}: which announces the completion of the goal. If the task is a QA problem, then the agent should submit the final answer in \textit{content} field; else the checking pipeline will automatically check the system status to judge the correctness.
\end{itemize}

\subsection{Prompt Example}
A prompt for OS evaluation consists of the instruction and the formulation of interaction trajectory.
An example of instruction prompt is: 

% \todo{update: PROMPT. please check @liuxiao} 

\lstset{
    backgroundcolor=\color[RGB]{245,245,244},
    breaklines=true,
    basicstyle=\ttfamily\small
}\begin{lstlisting}
You are an assistant that will act like a person, I'will play the role of linux(ubuntu) operating system. Your goal is to implement the operations required by me or answer to the question proposed by me. For each of your turn, you should first think what you should do, and then take exact one of the three actions: "bash", "finish" or "answer". 

1. If you think you should execute some bash code, take bash action, and you should print like this:

Think: put your thought here.

Act: bash

```bash
# put your bash code here
```

2. If you think you have finished the task, take finish action, and you should print like this:

Think: put your thought here.

Act: finish

3. If you think you have got the answer to the question, take answer action, and you should print like this:

Think: put your thought here.

Act: answer(Your answer to the question should be put in this pair of parentheses)

If the output is too long, I will truncate it. The truncated output is not complete. You have to deal with the truncating problem by yourself. Attention, your bash code should not contain any input operation. Once again, you should take only exact one of the three actions in each turn.
\end{lstlisting}

The trajectory is organized in CoT styles, and we use an 1-shot example to make model better understand the action space like the following.

\lstset{
    backgroundcolor=\color[RGB]{245,245,244},
    breaklines=true,
    basicstyle=\ttfamily\small
}\begin{lstlisting}
<USER>:
    {{ INSTRUCTION }}
    
    Now, my problem is:
    
    Tell me how many files are in the directory "/etc"?

<AGENT>:
    Think: To count the files in /etc, I need to print all the files in it.

    Act: bash
    
    ```bash
    ls /etc
    ```

<USER>:
    The output of the OS:
    cpi cron.hourly fuse.conf iproute2 lvm networkd-dispatcher protocols selinux tmpfiles.d [truncated because the output is too long]

<AGENT>:
    Think: The output has been truncated because it is too long, so I need to count files by script directly.

    Act: bash
    
    ```bash
    ls -1 /etc | wc -l
    ```

<USER>:
    The output of the OS:
    220

<AGENT>:
    Think: Now I get the answer, it is 220.

    Act: answer(220)

<USER>:
    Now, I will start a new problem in a new OS. My problem is:

    {{ PROBLEM }}

<LOOP>:
    <AGENT>:
        {{ RESPONSE }}

    Swicth AGENT.Action:
        "bash":
            <User>:
                The output of the OS:
                {{ OUTPUT }}
        "answer" or "finish":
            Executing Checking Script
            Exit Loop
        otherwise:
            Exit Loop
\end{lstlisting}
